\section{The Temptation of Saint Anthony}

\begin{quotex}
The right to be free to seek out the truth [involves] the duty to devote oneself to an ever deeper and wider search for it. \flright{\textsc{Pope John XXIII}, \emph{Pacem in Terris}}

\end{quotex} 

It is disheartening to listen to paid pundits — often earning 6 or even 7 digit salaries (in dollars) — making the most elementary errors in grammar, facts, and logic. In the USA, the rights to free speech and a free press are held to be sacrosanct. Unfortunately, this is usually understood in a nihilistic way, so that speech is an end in itself and not a means. However, as part of the US Constitution\footnote{\url{http://constitutionus.com/}} these rights do have a purpose, viz., those explicitly expressed in the preamble:

\begin{quotex}
We the People of the United States, in Order to form a more perfect Union, establish Justice, insure domestic Tranquility, provide for the common defence, promote the general Welfare, and secure the Blessings of Liberty to ourselves and our Posterity.

\end{quotex}
Since the Truth sets you free, the Blessings of Liberty will depend on speaking the truth. Unfortunately, the Truth is often hard to discover, so the citizen is obliged to devote himself to the search. At the least, this requires a good foundation in grammar, logic, and rhetoric.

I listened to a few minutes of a political talk show this morning. The Host, who is more rational than most, was spending all his time in explaining a common logical fallacy that has been making the rounds. Obviously, wasting time on elementary logic is not conducive to a deep search for truth, so I had to tune to the opera channel instead.

This is not the place and time for a detailed discussion of logical fallacies, but a quick survey of commonly believed misleading thoughts may benefit someone.

\paragraph{Aging Gracefully}
The supermarket tabloid covers often show headlines about post-menopausal models and superstars posing in their bikinis or boasting of their sex lives. Note that men don't purchase those publications.

There is a TV commercial about some arthritic old guy who manages to complete a footrace. That accomplishment is supposed to elicit our admiration for the human spirit. As an alternative, let me point out that \textbf{Ernst Junger} published \emph{Eumeswil}\footnote{\url{https://gornahoor.net/?p=9186}} when he was 82 years old. That is a real display of the human spirit.

So we propose this TV commercial. An old man, a former householder, becomes a sannyasi and achieves enlightenment.

\paragraph{Corporatism}
Left wingers misunderstand the term ``corporatism" as used in Fascist Italy to mean what is called a ``corporation" in the USA. Hence, they are incessantly complaining about ``corporate power". However, corporations, as for-profit business entities, are usually the engines of social leftism. They have even managed to overturn laws passed by elected officials in several states.

Actually, corporatism\footnote{\url{https://en.wikipedia.org/wiki/Corporatism}} has its roots in Ancient Greece, the Middle Ages, and was further developed in the 19th century by Pope Leo XIII. A modern economy can be divided into various business sectors. Ideally, each sector would involve the cooperation of government, business, and labour interests. Consider how trade, immigration, and social policies would differ if labour were part of the decision-making process.

In the current state of affairs, corporate power dominates life much more than the totalitarian governments of the past.

\paragraph{Statistical Intelligibility}
\textbf{Bernard Lonergan}, in \emph{Insight}, added the notion of statistical intelligibility to the classical notion of intelligibility. In brief,

\begin{quotex}
A statistical intelligibility is grasped by an inverse insight that there is no direct insight available. But while we often understand that many events cannot be functionally related to each other, we also may understand that an entire set of such events within a specific time and place will cluster about some average.

\end{quotex}
As a thought experiment, consider a naïve scientist who wants to study the tossing of a coin. He runs samples of 100 tosses, assigning 1 to heads and 0 to tails. Try as he might, he is unable to find any predictive pattern of 1's and 0's in his samples. At first, he is disappointed because a scientist's goal is ``direct insight", i.e., he wants to discover a law about the outcome of the tosses and how each toss is related to its predecessor.

He then graphs his samples, and surprisingly noticed a pattern in the graphs, forming what is now called the ``bell curve". He manipulates the numbers, eventually discovering the average number of heads in a sample and even derives secondary information such as standard deviations.

This is quite remarkable. Even though each individual coin toss is independent, and has no relation to what came before or after, as a whole, there is a recognizable pattern. This statistical intelligibility of the sample set is misunderstood in common discourse to be a judgment on the individual samples. This makes conversation quite difficult with those who don't understand statistical intelligibility.

\paragraph{Nominalism and Realism}
This hoary distinction is worth mentioning again. For those who may have missed it, nominalism is the notion that words are arbitrary human constructions; realism, then, is that words describe the essential nature of what they signify. There is a communications gulf between the two notions: a nominalist can make a statement in sincerity that sounds bizarre to a realist.

Consider this. Joe says, ``As an atheist, I believe in God." That is an obvious self-contradiction for the realist, since he understands the term `atheist' to entail the denial of God's existence. Now, a nominalist cannot logically reject that statement, since Joe's self-identity is his own choice. Nevertheless, that is a bit much. However, suppose Mary claims, ``I am a Mormon but I don't believe in God." A realist would say that Mary is not a ``real" Mormon, but merely a nominal Mormon. But nominalists are comfortable with that. Typically they might cite a study that says 20\% of Mormons are like Mary.

A farmer's product might be wheat; a builder's, a house; and so on. But for the intellectual class — e.g., journalists, lawyers, politicians — their very product is words. That is why, in the news media, concrete acts matter less than what people say, and beliefs are more important than truth. Hence, the discussion is primarily about who said what, what terms were used, how quickly did he say it, and the like. This class is entirely nominalist and has taken the place vacated by an ineffectual spiritual caste.

As \textbf{Martin Heidegger} said, language is the house of being in which humans dwell. Hence, whoever controls the use of language, then controls humans. That would be a blessing, if that class were dedicated to the truth and the moral betterment of the human race. You would think that someone whose life work depends on words, would be dedicated to the Trivium: grammar, logic, and rhetoric. Then language would reach a lofty place. Unfortunately, that is hardly the case, which Heidegger described:

\begin{quotex}
the widely and rapidly devastation of language not only undermines aesthetic and moral responsibility in every use of language; it arises from a threat to the essence of humanity.

\end{quotex}
The realist understands truth, goodness, and beauty to be realities to be discovered. The nominalist considers them to be social constructs. Hence, whoever controls discourse, thereby constructs society. In Heidegger's words:

\begin{quotex}
language surrenders itself to our mere willing and trafficking as an instrument of domination over beings.

\end{quotex}
\paragraph{Traditionalism vs Tradition}
A common misunderstanding, at least in certain limited circles, is that \textbf{Rene Guenon} inaugurated a new ``school" of thought called ``Traditionalism". Thus, they believe it be debated and critiqued like any other philosophical or intellectual system. Guenon is trying to be descriptive, not prescriptive. Hence, it makes no sense to become a ``traditionalist", but rather, one needs to follow a particular tradition. In his study of various traditions, he noticed certain common features. If you are so inclined, you don't refute him by an argument, but rather by denying those features. Ultimately, all the traditions accepted by Guenon as authentic are of two types:

\begin{itemize}
\item \textbf{The pagan traditions}, which derive from Hinduism. This includes Taoism, Buddhism, and the various folk religions of Europe. 
\item \textbf{The Abrahamic traditions}, which include Judaism, Islam, and Christianity. It is not always so neat, since the latter borrows heavily from the Greek and Roman religions. 
\end{itemize}
These are the main characteristics that Guenon identifies. I won't justify them all here, since there are plenty of posts that do so.

\begin{itemize}
\item Each tradition claims to be salvific. Other traditions may deny that claim, or claim exclusivity, but that is besides the point. 
\item There is an exoteric aspect to the tradition, suitable for most people. 
\item There is an esoteric aspect, which is open to a limited number. 
\item The traditions assume a common metaphysical framework. For example, in the Medieval Period, pagan philosophers (Plato, Aristotle, Plotinus), Jewish philosophers (Maimonides), Islamic philosophers (Averroes), and Christian philosophers (Augustine, Thomas Aquinas) were able to engage in fruitful exchanges of ideas. 
\item Social organization is based on caste, although their manifestations differed. This is not the ideal, but rather the actual, state of things since castes are based on human nature. Some types are more suited for spiritual and intellectual tasks, others for political and military tasks, and others for trade. A disordered society has not overcome castes, but rather has placed people where they don't belong. 
\end{itemize}
So a man who is familiar with Guenon's ideas and follows a Tradition such as Traditional Catholicism is ipso fact a ``Traditionalist". You can perhaps dispute about his level of understanding if you are so inclined, but then you will be judged by the same standard. Obviously, not everyone will follow an esoteric path or be adept at metaphysical ideas, but that is what we would expect. There is a statistical intelligibility even in this regard. Moreover, someone of the political or merchant caste would be unsuited for deep esoteric ideas.

\paragraph{Greek Popes}
There were a couple of centuries during the early Middle Ages when the Byzantines dominated the Roman papacy\footnote{\url{https://en.wikipedia.org/wiki/Byzantine_Papacy}}. Greek, Sicilian, and Syrians accounted for a large number of popes. You would think that if the Byzantines had any objections to the idea of the papacy, that would have been the time to bring it up.

On the theory of ``better late than never", the Greeks seem obsessed with it now. There are several Orthodox denizens on the Internet, particularly converts, who have a similar mentality to the ``Social Justice Warrior". They persistently battle Rome, or often its egregore, pointing out its alleged flaws. The somehow manage to find motes that few can see. They take pride in not having a pope, although logically, they should then reject the Patriarchs as the centre of spiritual authority.

It is well known that Guenon determined that the Roman religion had lost some of her traditional elements. Of course, that is not a metaphysical certainty, rather a contingent fact that may or may not be true. Then, again, that is to be expected according to traditional teachings. Several saints have had premonitions about the failure of the popes and priests to uphold the true tradition.

Some Orthodox Warriors believe they are immune to that decline. They complain that Guenon was ignorant about the East, or he would not have had such a view about the state of the Christian religion. That is ridiculous. From Rome's perspective, which was also Guenon's, there was a schism between East and West. A schism is solely an administrative dispute about which priests are in charge of things. There is no difference between the religion, or tradition, of the East and West. There is at most a difference in emphasis, or in customs that may appear alien.

I recently encountered a blog that was critical of Catholic saints. For example, the author was critical of St. John of the Cross, whose three stages of the spiritual life can be traced back to Dionysius. He then turned to St. Ignatius of Loyola, even quoting from the Philokalia to condemn the Jesuit founder. This irony was lost on that poor fellow. The compiler of the Philokalia, St. Nicodemos, was actually quite impressed by St. Ignatius, even translating his \emph{Spiritual Exercises} into Greek! You are taking a big risk to your spiritual life when you engage with these Orthodox warriors.

It is a shame to see self-loathing Westerners reject the Traditions of their fathers and apostatize from their faith. There is much they have to give up, and it is seldom justifiable. The only case is a state of emergency, which occurs when there is no reliable Roman priest available. But, in that case, you will not become a faux Warrior. The irony is that while Western Lutherans are converting to Orthodoxy, Russian intellectuals are converting to Lutheranism\footnote{\url{http://www.patheos.com/blogs/geneveith/2017/08/lutheranism-appeals-russians/}}. I suppose that is not so surprising, since Fr. Seraphim Rose often boasts about how anti-intellectual Orthodoxy is.

There must be something about the Greek view of things that values Strife. Years ago, when the Greek church was built in Boca Raton, it attracted many Russians who did not have a church of their own. On one occasion, there was a dispute about the appointment of a new priest that became so serious, so that the Police had to intervene during Mass to break up the fight. Since then, they have managed to produce the best Greek Festival from miles around.

At the end of the Latin mass, we pray for reconciliation with the East. The Orthodox warriors, on the other hand, are often addicted to strife, dissension, and the party spirit\footnote{\url{https://biblehub.com/context/galatians/5-19.htm}}. The best course would be to end the schism. Perhaps, then, the Russian patriarch would become the next pope, which would be a blessing for all. Then we would just need the Emperor.

\paragraph{Conservatism vs Populism}
Finally getting around to spring cleaning, I found some old magazines\footnote{\url{https://www.chroniclesmagazine.org/}} from around 20 years ago. One article listed the following points as fundamental to conservatism as it was understood at that time:

\begin{itemize}
\item Immigration control (as opposed to open borders) 
\item Economic nationalism (as opposed to globalism) 
\item Military isolationism (as opposed to adventurism) 
\item Social issues (not just sexual) 
\end{itemize}
Curiously, those same points are dismissed as ``populism" and soi-disant conservatives today typically reject those points.

\paragraph{The Descent into Hell}
There have been several movies and books describing wonderful and peaceful postmortem experiences in Heaven. There is some truth about the experiences, although it only signifies the detachment of the soul from the burden of the body. Of course, that seems peaceful, but the experience can be replicated in deep meditation, so it is not a sign of death and a return.

In \textbf{Dante}'s \emph{Divine Comedy}, we learn that the path to Heaven passes through Hell. The fear and dread of traversing Hell prevents most people from even beginning the journey.

\paragraph{The Problem of Masturbation}
Last night, I was awakened in the middle of the night by a bad dream. I had a little erection so I considered masturbating in order to relax and get back to sleep. I couldn't quite remember the last time I did that, so perhaps it might be OK to try it again. It all sounded so logical. However, further thought — for several reasons — dissuaded me from acting on that thought.

We can speak openly of such things today. This is unlike \textbf{Athanasius}'s \emph{Life of Saint Anthony} which simply described the attack of demons during the night, without describing the temptations in detail. Obviously, that act creates an opening for demonic influences to enter. Demons don't typically announce themselves as evil.

Masturbation is one of the foundations of the modern world along with sterility, death, and drugs. They all make the population stupider and weaker. Boys will pay girls to play online games with them, or send tips to sexy models on Instagram. I know a computer engineer making a high six digit salary running a network of webcam girls. The power of that temptation is nearly overpowering.



\flrightit{Posted on 2017-08-17 by Cologero }

\begin{center}* * *\end{center}

\begin{footnotesize}\begin{sffamily}



\texttt{Fred on 2017-08-18 at 03:01 said: }

An Orthodox critcising St. John of the Cross only shows prejudiced ignorance and hostility. There is an important athonite saint who read the dark night of the soul and was very surprised to discover noetic prayer in western christianity. He called St.John ``a western hesychast who use a different language from the Fathers".

If people with actual experience of contemplation believe this why sould we bother with the opinions of those who had no experience of prayer and only judge according to futile mental rationalizations?

Also I would be very wary to say the pagan traditions derive from Hinduism.

Other indo-european pagan traditions lack what hinduism borrowed from the people they conquered (Rudraism and tantra were not beliefs of the aryans).

Taoism is the native religion of China and is obviously not an indoeuropean religion.


\hfill

\texttt{Fred on 2017-08-18 at 03:08 said: }

A paradox: sometimes I think the right Church between the Catholic and the Orthodox is the one who will be humble enough to accept the claims of the other.

(I had to make a double post because the button to send the message disappear atr the bottom of the page after I write a few lines, I'll take this as a sign not to be verbose)


\hfill

\texttt{Cologero on 2017-08-18 at 03:40 said: }

Thanks, Fred, for the anecdote about St. John of the Cross.

Yes, ``derive from" is too strong a phrase, but that does not mean there is not a relationship. I had in mind the Vedas, not any later additions to Hinduism. Based on the works of Dumezil, there is a homomorphism among the ancient pagan mythologies. Check out \textit{Interpretatio graeca}\footnote{\url{https://en.wikipedia.org/wiki/Interpretatio_graeca}}: even the ancient Greeks, Romans, and Germanics noted the similarity. They simply were not familiar with Vedic gods. So the idea of a proto-indo-european religion\footnote{\url{https://en.wikipedia.org/wiki/Proto-Indo-European_religion}} is not so far-fetched.

In \emph{Shakti and Shakt}, Arthur Avalon demonstrates that Taoism is a derivative of the Upanishads. I defer to his expertise in this opinion.


\hfill

\texttt{jean lave on 2017-08-21 at 14:47 said: }

Hello

Taoism is esoterism.

Confucianism is exoterism.

Both together are the native religion of China.

According to R.Guenon St.John of the Cross was a mystic,and mysticism is different from initiation.

Hesychast is esoterism of orthodox christianity.

Beside mysticism is specific to christianity.(doesn't exist in orient)

Shall we admit that Guenon was one of ``those who had no experience of prayer and only judge according to futile mental rationalizations?"

Sorry for my english which is not my mother tongue.

thank you for your blog.

in the Name of HE by whom everything exists.

PEACE


\hfill

\texttt{Cologero on 2017-08-21 at 21:40 said: }

M. Lave,

It was a French priest, Father P. L. Wieger, S.J, who noted that Taoism is likely a development from the Upanishads. Please look at \textit{Shakti and Shakta}\footnote{\url{http://www.mysticknowledge.org/Shakti_and_Shakta-By_Arthur_Avalon.pdf}} by Arthur Avalon. China had a native folk, shamanic religion long before Lao Tzu and Confucius. 

I can't judge about R. Guenon, as we know nothing about his experiences in Egypt. That is why I don't agree with his ``impersonal" approach. Instead, it is possible to be ``personal" in an impersonal way. For example, I believe it is helpful to describe in some detail how to deal with tempting thoughts and images, or how to achieve concentration without effort. That is the bare beginning of the hesychast path, in case you don't have your own hesychast teacher. In that case, it is better to at least discover what you are capable of. Hesychasts won't bother to read this.

I'll accept Fred's comment on John of the Cross. How can the unitiated judge an initiate? or even a ``mystic" for that matter. If John of the Cross achieved theosis, then he is a hesychast.


\hfill

\texttt{Fred on 2017-08-23 at 13:23 said: }

I am familiar with the writings of Guenon and can safely affirm that he fell into a deep misconception by calling St. John of the cross a simple non esoteric mystic, as if St.John was just one of those pious old women who had visions (the actual mystics he had in mind when he used that word)

And yes Guenon never followed a path of Christian esotericism, so on this issue (prayer and hesychasm) his judgment on st. John of the cross is weaker compared to that of st. Silouan (the saint I mentioned in my previous comment, although it was expressed through his disciple Sophrony). 

For a deeper explanation of why the typical Sufi/Sharia distinction of Guenon is applied badly to Christianity read Jean Borella.

I long struggled with this problem and concluded the ipse dixit of Guenon is not enough to solve it.


\hfill


\end{sffamily}\end{footnotesize}
